% A made-up manuscript for the demo project of Papers Review (its screenshots): no real study, no real authors.
\documentclass[11pt]{article}
\usepackage[a4paper,margin=2.8cm]{geometry}
\usepackage{lineno,amsmath,lipsum}
\usepackage[T1]{fontenc}
\usepackage{lmodern}
\linenumbers
\title{Hysteresis of a marine ice-sheet margin in a two-stage toy model}
\author{A. N. Author$^{1}$ and B. Coauthor$^{2}$\\[2pt]
\small $^1$Institute of Imaginary Glaciology \quad $^2$Department of Made-up Sciences}
\date{Preprint, demo manuscript}
\begin{document}
\maketitle
\begin{abstract}
Marine ice sheets resting on bedrock that deepens inland can retreat irreversibly once their grounding line passes a
threshold. We present a two-stage toy model of a marine margin that couples a flux condition at the grounding line to
a slowly adjusting ocean forcing. The model reproduces the hysteresis loop of more complex flowline models at a small
fraction of their cost, and shows that the width of the loop depends mainly on the basal friction exponent. Under a
warming of 1.5\,K the margin crosses its upper tipping point within two centuries, while recovery requires a cooling
well below the pre-industrial state.
\end{abstract}

\section{Introduction}
The stability of marine ice-sheet margins has been debated since the first theoretical treatments showed that a
grounding line on a retrograde bed has no stable position. Later work derived a flux condition at the grounding line
that made the argument quantitative, and flowline models have since mapped the hysteresis of idealized margins in
detail. These models remain expensive to run over the many thousands of years and the wide range of forcing that an
assessment of irreversibility requires.

Here we take the opposite approach. We reduce the margin to two stages, the grounded sheet and a shelf, and keep only
the processes that set the position of the grounding line. The aim is not to replace the flowline models but to make
their main result, the width of the hysteresis loop, easy to explore. We ask three questions: how wide is the loop,
what sets its width, and how fast is it crossed under a realistic warming?

\section{Model}
The grounded stage has a thickness $H$ and a length $L$. Mass is gained by accumulation $a$ over the length and lost
through the flux $q_g$ at the grounding line, so that
\begin{equation}
  \frac{d}{dt}(H L) = a L - q_g .
\end{equation}
Following the flux condition of the boundary-layer theory, the flux at the grounding line scales with the ice
thickness there to a high power,
\begin{equation}
  q_g = C\, h_g^{\,(m+n+3)/(m+1)} ,
\end{equation}
where $m$ is the basal friction exponent, $n=3$ the exponent of the flow law and $C$ a constant that collects the
material and friction parameters. The ocean forcing enters as a melt rate at the base of the shelf that relaxes
towards a target set by the temperature anomaly $\Delta T$ with a time scale $\tau$ of a few decades.

\section{Results}
The steady states of the model form a fold in the plane of forcing and margin position. Between the two folds, three
steady states coexist for the same forcing, of which the middle one is unstable. The loop is widest for a linear
friction law, $m=1$, and narrows as $m$ decreases towards a plastic bed. For $m=1/3$ the upper tipping point lies at
a warming of about 1.2\,K, and the margin only regrows once the ocean has cooled by 0.8\,K below its initial state.

The time to cross the loop depends on the rate of warming. Under a warming of 1.5\,K reached within a century, the
grounding line passes the upper fold after 120 to 190 years, depending on the friction law, and the retreat then
proceeds at a rate set by the accumulation alone. The transient lag between forcing and response is therefore
dominated by the slow mass balance of the grounded stage and not by the ocean.

\section{Discussion}
\lipsum[1-4]

\section{Conclusions}
\lipsum[5-6]
\end{document}
